\section{Fazit}

\subsection{Zusammenfassung}
\begin{itemize}
  \item \textbf{Mercator erfüllt das Ziel einer analytischen Datenanwendung für Insider-Trade-Analyse.} % Repo-Evidenz prüfen
  \item Die Anwendung verbindet öffentliche API-Datenquellen, Validierung, lokale Filterung, Datenenhältung und interaktive Explorationstools in einer kohärenten Pipeline. % Repo-Evidenz prüfen
  \item Die Raw/Clean-Trennung ist ein zentraler Architekturbeitrag, der Audit-Anforderungen, Debugging und Reprodzierbarkeit unterstützt. % Eigene methodische Ableitung
  \item Die Kombination aus MongoDB (flexibel) und MySQL (strukuriert) ist fachlich begründbar und performant. % Eigene methodische Ableitung
  \item Lokale Gate-Filter reduzieren API-Aufrufe deutlich, ohne Funktionalität zu opfern. % Interne Projektspezifikation
  \item Streamlit ermöglicht schnelle und intuitive explorative Analyse ohne Browser-Framework-Overhead. % Repo-Evidenz prüfen
\end{itemize}

\subsection{Erfüllung der Anforderungen}
\begin{itemize}
  \item \textbf{Modul-Anforderungen (Datenbanken 2):}
    \begin{itemize}
      \item Öffentliche Datenquelle (FMP API) ✓
      \item Pandas-Verarbeitung ✓
      \item MySQL-Speicherung ✓
      \item MongoDB-Speicherung ✓
      \item Interaktive Streamlit-UI ✓
      \item Visualisierungen (Charts, KPIs, Tabellen) ✓
    \end{itemize}
  \item \textbf{Akademische Qualität:}
    \begin{itemize}
      \item Methodologie-Dokumentation ✓
      \item Nachvollziehbare Architektur ✓
      \item Transparente Gate- und Scoring-Logik ✓
      \item Kritische Reflexion der Grenzen ✓
    \end{itemize}
\end{itemize}

\subsection{Technische und methodische Beiträge}
\begin{itemize}
  \item \textbf{Raw/Clean-Architektur:} Zeigt, wie polyglot persistence praktisch in einer analytischen Anwendung umgesetzt wird. % Eigene methodische Ableitung
  \item \textbf{API-schonende Verarbeitung:} Gate-Filter vor teureren API-Calls reduzieren Kosten und Rate-Limit-Probleme. % Interne Projektspezifikation
  \item \textbf{Standardisierte Trade-Akkumulation:} 3-Tage-Fenster mit fachlicher Gruppierung reduziert Rauschen. % Interne Projektspezifikation
  \item \textbf{Transparente Scoring-Heuristik:} Regelbasierter Score ohne Black-Box-ML. % Eigene methodische Ableitung
  \item \textbf{Barrierefreie UI-Gestaltung:} Text-basierte Status-Labels neben Farben (WCAG-Konformität). % Repo-Evidenz prüfen
\end{itemize}

\subsection{Realistische Ausbaustufen}
\begin{itemize}
  \item \textbf{Kurzfristig (0--3 Monate):}
    \begin{itemize}
      \item Monitoring und Logging ausbauen.
      \item Test-Coverage erhöhen (Unit / Integration / E2E). % Repo-Evidenz prüfen
      \item Fehlerbehandlung robuster machen (Circuit-Breaker, Retry-Strategien).
      \item UI-Performance optimieren (Caching, Query-Optimierungen).
    \end{itemize}
  \item \textbf{Mittelfristig (3--12 Monate):}
    \begin{itemize}
      \item Zeitreihen- und Trend-Analysen erweitern.
      \item Backtesting-Framework für Score-Validierung.
      \item Netzwerk-Analysen (Insider $\leftrightarrow$ Company $\leftrightarrow$ Sektor).
      \item Redis-Cache fur API-Payloads und technische Metriken.
      \item Parallele API-Aufrufe (asyncio).
    \end{itemize}
  \item \textbf{Mittellangfristig (12+ Monate):}
    \begin{itemize}
      \item Machine Learning zur Score-Kalibrierung (bei ausreichend historischen Daten).
      \item Mehrsprachige UI (Internationalisierung).
      \item Mobile App (React Native oder Flutter).
      \item Horizontale Skalierbarkeit (Kubernetes, Load-Balancing).
      \item Multi-Tenant-Unterstützung (SaaS-Modell).
    \end{itemize}
\end{itemize}

\subsection{Ausblick und Nachhaltigkeit}
\begin{itemize}
  \item \textbf{Code-Qualität:} Das Projekt folgt best practices (Logging, Typisierung, Fehlerbehandlung) und ist wartbar. % Repo-Evidenz prüfen
  \item \textbf{Dokumentation:} Methodologie, API-Spezifikation, README sind ausführlich vorhanden. % Repo-Evidenz prüfen
  \item \textbf{akademische Weiternutzung:} Nachfolge-Projekte können auf dieser Basis aufbauen und vertiefen (z. B. Statistik, NLP, Advanced Analytics). % Eigene methodische Ableitung
  \item \textbf{Offene Fragen für Future Work:}
    \begin{itemize}
      \item Können Insider-Käufe unter bestimmten marktlichen Bedingungen als Signal für zukünftige Überperformance genutzt werden? (Empirische Validierung erforderlich.) % Externe Quelle nötig
      \item Wie können Szenarien-Analysen (z. B. Sektoren, Marktphasen) in das Scoring integriert werden? % Externe Quelle nötig
      \item Welche alternativen API-Quellen (z. B. Alternative Data, Sentiment) könnten die Signalqualität verbessern? % Externe Quelle nötig
    \end{itemize}
\end{itemize}
